%% PLANTILLA PARA ARTICULOS

\documentclass[10pt,twoside,spanish,a4paper]{article}
%%%%%%%%%%%%%%%%%%%%%%%%%%%%% no tocar
% Alfredo Cuesta. 01/01/2013. acuestai@ucm.es

%si se quiere usar ARIAL hay que poner estas 2 lineas:
%\usepackage[scaled]{uarial} %DESCOMENTAR
\renewcommand*\familydefault{\sfdefault}
%

\usepackage[compact]{titlesec}
 \titlespacing\section{0pt}{12pt plus 0pt minus 0pt}{0pt plus 0pt minus 0pt}
 \titlespacing\subsection{0pt}{5pt plus 0pt minus 0pt}{0pt plus 0pt minus 0pt}
 \titlespacing\subsubsection{0pt}{0pt plus 0pt minus 0pt}{0pt plus 0pt minus 0pt}


%\usepackage[T1]{fontenc}
\usepackage[utf8]{inputenc}
%\usepackage[latin9]{inputenc}
\usepackage{
amssymb,
graphicx,
color,
url
%wrapfig %Inclusión de gr·ficos al lado de texto
}

\linespread{1.3}%separacion de lineas
\parindent=0.0in%indentancion de parrafo
\parskip=6pt%separacion entre el parrafo siguiente
\setlength{\topmargin}{-0.5in}
\setlength{\oddsidemargin}{0in}
\setlength{\evensidemargin}{0in}
\setlength{\textwidth}{16cm}%{4.72in}
\setlength{\textheight}{22cm}%{7.08in}

\usepackage{babel}
\addto\shorthandsspanish{\spanishdeactivate{~<>}}
% correct bad hyphenation here
\hyphenation{nu-es-tro co-no-ci-mien-to par-ti-cu-lar ca-rac-te-ris-ti-cas ni-vel se-gui-mien-to que-re-mos ca-ta-lo-go}
\date{}


% --- customization ---

\renewcommand\labelitemi{$\bullet$}
\let\olditemize\itemize
\renewcommand{\itemize}{
  \olditemize
  \vspace{-0pt}
  \setlength{\itemsep}{6pt}
  \setlength{\parskip}{0pt}
  \setlength{\parsep}{0pt}
  %\setlength{\topsep}{-15pt}
}

\let\oldenumerate\enumerate
\renewcommand{\enumerate}{
  \oldenumerate
  \vspace{-2pt}
  \setlength{\itemsep}{1pt}
  \setlength{\parskip}{0pt}
  \setlength{\parsep}{0pt}
  %\setlength{\topsep}{-15pt}
}

\usepackage{verbatim}
\usepackage{fancyhdr}
\usepackage{multirow}
\usepackage{colortbl}
\usepackage{verbatim}
\usepackage{amsmath}

\definecolor{olivegreen}{rgb}{0.0,0.5,0.0}

\begin{document}

\renewcommand{\tablename}{Tabla}
\renewcommand{\figurename}{Figura}
\title{
  \vspace{-0.2cm}
	\begin{center}
		\large
		\begin{tabular}{rrl}
		\multirow{4}{*}{\includegraphics[width=1.7cm]{UCM.eps} \hspace{1cm}}&Grado: & Ingeniería Electrónica de Comunicaciones\\
		&Asignatura:& Procesamiento de Señales \\
		&Profesores:  & Jesús Chacón Sombría \\
		\end{tabular}
	\end{center}
	\rule{\textwidth}{0.4pt}
	%
	\vspace{-1.5cm}
	\begin{center}
		\color{blue} \Huge
		\begin{tabular}{c}
	        \textbf{Ejercicios}\\
		\textbf{Introducción: señales y sistemas}\\
		\end{tabular}
	\end{center}
	\vspace{-0.5cm}
	\rule{\textwidth}{0.4pt}
}
\maketitle
 \thispagestyle{empty}
\pagestyle{fancy}

\lfoot{\includegraphics[width=0.8cm]{UCM.eps}}
\cfoot{\vspace{-0.6cm} Ingeniería Electrónica de Comunicaciones }
\rfoot{\vspace{-0.6cm} \thepage}
\lhead{Ejercicios}
\rhead{Muestreo en el dominio del tiempo}
\renewcommand{\headrule}{ \vspace{-0.2cm} \rule{\textwidth}{0.4pt}}
\renewcommand{\footrule}{\vspace{0.5cm} \rule{\textwidth}{0.4pt}\\ \vspace{-0.3cm} \\}

%**************************************************************<- para indicar pregunta de la pr·ctica
\newcounter{ctPreg}
\setcounter{ctPreg}{1}
\newcommand{\itemP}{ \item[{\color{blue}\fbox{\textbf{Ejercicio \arabic{ctPreg}:}}}] \addtocounter{ctPreg}{1}}  %
\newcommand{\itemB}{ \item[{\color{blue}$\bullet$} ]    }
\newcommand{\itemR}{ \item[{\color{blue}$\triangleright$} ]    }       %
%**************************************************************

\vspace{-8em}

	\begin{itemize}
		\itemP Caracterizar las siguientes señales según su \emph{dimensión}, y continuidad de las variables dependientes e independientes:
		\begin{itemize}
				\itemB $x(t)=\begin{pmatrix}4t \\ e^{j 2t}\end{pmatrix}$
				\itemB $x(n)=4 \textrm{e}^{j\frac{3}{5}n}$
				\itemB $x[n]=3 \textrm{e}^{j\frac{2}{\pi}n}$
				\itemB $x[n]=\lfloor sen (\frac{2\pi n}{20})\rfloor$
				\itemB $x(t,x)=\textrm{sen}(\omega t-kx)$
		\end{itemize}
		
		\itemP Clasifica las siguientes señales según su \emph{periodicidad} y su \emph{paridad}. Para las señales periódicas, indica su periodo.
		\begin{itemize}
				\itemB $x(t)=4 \textrm{e}^{j\frac{3}{5}t}$
				\itemB $x(t)= 3 t$
				\itemB $x(t)= 3 t^2$
				\itemB $x[n]= \textrm{sen}(3 n)$
				\itemB $x[n]= \textrm{cos}(\frac{\pi n}{20})$
				\itemB $x[n]= \textrm{cos}(3 n - 4),0\leq n \leq N_0$ 
		\end{itemize}
		
		\itemP Clasifica las siguientes señales según su \emph{energía} y \emph{potencia} (finita o infinita).
		\begin{itemize}
				\itemB $x(t)=t$
				\itemB $x(t)=u(t)$ (escalón unitario)
				\itemB $x(t)=\textrm{e}^{- j 4 t}$
				\itemB $x(t)=4t,0\leq t \leq 10$
		\end{itemize}

		\clearpage
		\itemP Indica si las siguientes señales son reales o complejas.
		\begin{itemize}
				\itemB $x(t)=e^{-3 t}$
				\itemB $x(t)=e^{j 2\pi t}$
				\itemB $x[n]=e^{j \pi n}$
				\itemB $x(t)=\textrm{cos}(4t) + j \textrm{sen}(5t)$
				\itemB $x(t)=e ^{j \pi} e^{-j 3t}$
		\end{itemize}
		
		\itemP Realiza las operaciones indicadas sobre las siguientes señales $x_1(t)=\textrm{sin(t)}$, $x_2(t)=\textrm{cos}(\frac{2\pi}{5} t)$, $x_3[n]=[0,3,1,1]$, $x_4[n]=\textrm{e}^{j 4 n}$.
		\begin{itemize}
				\itemB $x_1(t) * x_2(t)$
				\itemB $x_1(t) + x_2(t)$
				\itemB $(x_1*x_2)(t)$
				\itemB $x_3[n] / x_4[n]$
				\itemB $x_3[n] - x_4[n]$
				\itemB $(x_3*x_4)[n]$
		\end{itemize}
		
		\itemP Calcular la convolución discreta $(x*y)[n]$ de las siguientes señales:
		\begin{itemize}
				\itemB $x[n]=[1,4,2,0,1]$ e $y[n]=[4,2,3]$
				\itemB $x[n]=4\delta[n-6]$ e $y[n]=\textrm{e}^{j\frac{2\pi}{5} n}$
				\itemB $x[n]=u[n]$ (escalón) e $y[n]=a^n$
		\end{itemize}
		
		\itemP Calcular la expansión en serie de Fourier de las siguientes señales:
		\begin{itemize}
				\itemB Onda cuadrada
				\itemB Onda triangular
		\end{itemize}
		
		\itemP Calcular la transformada de Fourier de las siguientes señales, para el caso continuo y el discreto.
		\begin{itemize}
				\itemB $\delta(t-t_0),\delta[n-n_0]$ (Impulso)
				\itemB $\Pi(\frac{t}{T}),\Pi[\frac{n}{N}]$ (Pulso cuadrado causal y no causal).
				\itemB $\wedge(t),\wedge[n]$ (Pulso triangular (duración T) causal y no causal)
		\end{itemize}
		
		\itemP Determinar si la señal es periódica en $n$ y en $\omega$.
			
		\begin{itemize}
				\itemB $x[n]=2\textrm{e}^{j(2\pi/8)n}$
				\itemB $x[n]=5\textrm{e}^{j(3/4)n}$
				\itemB $x[n]=cos(2\pi n)$
				\itemB $x[n]=\frac{1}{n}$
		\end{itemize}
		
		\itemP Determinar la estabilidad de los siguientes sistemas continuos:
		\begin{itemize}
				\itemB $H(s)=\frac{1}{s}$
				\itemB $H(s)=\frac{s+2}{s+4}$
				\itemB $H(s)=\frac{1}{s-3}$
				\itemB $H(s)=\frac{s+1}{s^2+25}$
		\end{itemize}
		
		\itemP Determinar la estabilidad de los siguientes sistemas discretos:
		\begin{itemize}
				\itemB $H(z)=\frac{z}{z-1.5}$
				\itemB $H(z)=\frac{z}{(z-1)^2}$
				\itemB $H(z)=\frac{z}{z-0.9}$
				\itemB $H(z)=\frac{z}{z^2+1}$
		\end{itemize}
		
		\itemP Determinar en cada caso si la respuesta del sistema a la señal de entrada es o no acotada.
		\begin{itemize}
				\itemB $H(s)=\frac{s}{s+4}$, entrada: $u(t)=\textrm{sin}(4t)$
				\itemB $H(s)=\frac{s^2}{s^2+16}$, entrada: $u(t)=\textrm{cos}(2t) + \textrm{sin}(4t)$
				\itemB $H(z)=\frac{z~\textrm{sin}(\frac{2\pi}{7})}{z^2-2z~\textrm{cos}(\frac{2\pi}{7})+1}$, entrada: $u[n]=3\textrm{cos}({-\frac{4\pi}{14}n})$
				\itemB $H(z)=\frac{z}{z^2+2}$, entrada: $u[n]=0.9^n$
		\end{itemize}
		
		\itemP Obtener, utilizando la transformada de Fourier y Laplace, la respuesta temporal de los siguientes sistemas a las correspondientes señales de entrada (se asume $u(t)=0, \forall t<0$).
		\begin{itemize}
				\itemB $h(t)=\textrm{e}^{-2t}$, entrada: $u(t)=4~\textrm{cos}(3t)+10~\textrm{sin}(5t-15)$.
				\itemB $h(t)=\textrm{sin}(4t)$, entrada: $U(s)=\frac{s}{s^2+5}$
				\itemB $H(s)=\frac{s+3}{s+10}$, entrada: $u(t)=t^2$
		\end{itemize}
		
		\itemP Obtener, utilizando la transformada de Fourier y Z, la respuesta temporal de los siguientes sistemas a las correspondientes señales de entrada.
		\begin{itemize}
				\itemB $H(z)=\frac{z}{z+0.8}$, entrada: $u[n]=0.5^n + 4\cdot0.8^{n-1}$.
				\itemB $H(z)=\frac{z}{z-0.8}$, entrada: $u[n]=n$
				\itemB $H(z)=\frac{z^2}{z^2-0.3z+0.4}$, entrada: $u[n]=\textrm{cos}(\frac{\pi}{6}n)$
		\end{itemize}
		
		\itemP Obtener, utilizando la transformada de Fourier, la respuesta permanente de los siguientes sistemas a las correspondientes señales de entrada (suponer que las señales de entrada son causales).
		\begin{itemize}
				\itemB Sistema: $\frac{1}{s+4}$, entrada: $x(t)$ onda cuadrada con $T=2$.
				\itemB Sistema: $\frac{z}{z-0.8}$, entrada: $2\textrm{cos}(\frac{\pi}{6}n)+\textrm{sin}(\frac{\pi}{2}n)$.
				\itemB Sistema: $\frac{z}{z-1.2}$, entrada: $\textrm{cos}(\frac{2\pi}{10}n)$.
		\end{itemize}
	\end{itemize}
	
\end{document}
